\documentclass[11pt,letterpaper]{article}
\usepackage[T1]{fontenc}
\usepackage{lmodern,microtype}
\usepackage[margin=0.85in]{geometry}
\usepackage{amsmath,amssymb,amsthm,mathtools}
\usepackage{booktabs,longtable,array}
\usepackage{xcolor,enumitem,fancyhdr}
\usepackage{hyperref,xurl}
\hypersetup{colorlinks=true,linkcolor=blue!50!black,urlcolor=blue!50!black,citecolor=blue!50!black}
\setlength{\parindent}{0pt}
\setlength{\parskip}{5pt}
\setlength{\emergencystretch}{2em}
\setlist{nosep,leftmargin=1.5em}
\pagestyle{fancy}
\fancyhf{}
\lhead{\small Direct Diophantine consistency}
\rhead{\small Pooled powers and integral projection}
\cfoot{\thepage}
\setlength{\headheight}{14pt}
\newtheorem{theorem}{Theorem}[section]
\newtheorem{lemma}[theorem]{Lemma}
\newtheorem{proposition}[theorem]{Proposition}
\newtheorem{corollary}[theorem]{Corollary}
\theoremstyle{remark}
\newtheorem{remark}[theorem]{Remark}
\newcommand{\Z}{\mathbb Z}
\newcommand{\N}{\mathbb N}
\newcommand{\Con}{\operatorname{Con}}
\newcommand{\ZFC}{\mathrm{ZFC}}
\newcommand{\ZF}{\mathrm{ZF}}
\newcommand{\Ccal}{\mathrm{C222}}
\newcommand{\pp}[1]{\texttt{\detokenize{#1}}}
\newcommand{\MainTerms}{7,726}
\newcommand{\MainHeight}{45,161}
\newcommand{\MainRowHeight}{128}
\newcommand{\FewTerms}{15,158}
\newcommand{\FewHeight}{334,952}
\newcommand{\FewRowHeight}{384}

\title{\bfseries Pooling signed power certificates before integral projection\\[5pt]
\large A direct quartic consistency equation in 218 integer variables,\\
with a 216-variable companion}
\author{Prepared for Vladimir Reshetnikov and the ProveIt research program\\
\small AI-assisted research and independent internal audits}
\date{25 October 2026}
\begin{document}
\maketitle
\begin{abstract}
We give a fully specified polynomial \(P_{218}\in
\Z[x_0,\ldots,x_{217}]\) of total degree four for which an ordinary
argument formalizable in ZFC proves
\[
  \bigl(\forall \mathbf{x}\in\Z^{218}\;P_{218}(\mathbf{x})\ne0\bigr)
  \quad\longleftrightarrow\quad\Con(\ZFC).
\]
It is the sum of squares of 155 explicitly printed quadratic residuals.
Its collected expansion has \MainTerms\ nonzero monomials and maximum
absolute coefficient \MainHeight. The inspected 225-variable direct
equation has 162 residuals, 7,867 monomials, and height 128,446.
Thus the new principal equation improves all four size measures,
with degree unchanged. A second exact construction has 216 variables
and 153 residuals, at the measured cost of \FewTerms\ monomials and
height \FewHeight.

The construction combines the factorial-free arithmetic source of
ProveIt source 157 with the fixed-base power pooling of source 158,
then applies explicit integral projections and residual row operations.
The combination requires a new check of the signed bootstrap and of
every retained square-frame atom. Every constant of the final
polynomials is a literal integer. There is no hidden universal-polynomial
parameter or Turing-machine coefficient tuple.

The inherited C222/ZF/ZFC refutation calibration is identified exactly.
The new semantic composition and algebraic maps are ordinary proofs,
with independent finite coefficient checks. This is a concrete improvement
in the direct quartic branch, not a global simplicity minimum, a
consistency decision, or a newly compiled end-to-end formal proof.
\end{abstract}
\tableofcontents

\section{The delivered equation and the scope of the result}

All independent unknowns of the delivered equations range over
\emph{unrestricted integers}.
Let \(R_0,\ldots,R_{154}\) be the literal quadratic polynomials printed
in Appendix~\ref{app:residuals}. The principal equation is
\begin{equation}\label{eq:principal}
 \boxed{\displaystyle P_{218}(x_0,\ldots,x_{217})
       =\sum_{j=0}^{154}R_j(x_0,\ldots,x_{217})^2=0.}
\end{equation}
The appendix, together with this finite sum, completely defines one
polynomial. The machine-readable coefficient list and complete collected
expansion are supplied as well. The indices of the sum are fixed;
they are not additional unknowns.

\begin{theorem}[Consistency equation]\label{thm:main}
For the effective presentation of ZFC used by the repository's named
refutation calibration, ZFC proves
\[
 \exists\mathbf{x}\in\Z^{218}\; P_{218}(\mathbf{x})=0
 \quad\longleftrightarrow\quad
 \exists p\;\operatorname{Refutation}_{\ZFC}(p).
\]
Consequently ZFC proves that \eqref{eq:principal} has no integer
solution if and only if \(\Con(\ZFC)\).
The analogous statement holds for the supplied
\(P_{216}\in\Z[x_0,\ldots,x_{215}]\).
\end{theorem}

The theorem composes the arithmetic equivalences and integer polynomial
maps below with the inherited syntactic calibration. It asserts neither
side of the consistency equivalence by itself. Evaluating these
polynomials is not a method for deciding consistency.

\begin{table}[ht]
\centering\small
\begin{tabular}{lrrrrr}
\toprule
Equation & Variables & Rows & Degree & Terms & Height\\
\midrule
Source 158, \(P_{225}\) &225&162&4&7,867&128,446\\
New principal \(P_{218}\)&218&155&4&\MainTerms&\MainHeight\\
New companion \(P_{216}\)&216&153&4&\FewTerms&\FewHeight\\
\bottomrule
\end{tabular}
\caption{Exact counts over \(\Z\). Rows are quadratic residuals before
taking the sum of squares. Height is the maximum absolute
\emph{expanded polynomial} coefficient.}
\label{tab:comparison}
\end{table}

The principal residual height is \MainRowHeight, and all 218 declared
coordinates occur. The companion residual height is \FewRowHeight.
The companion uses two fewer coordinates but is less sparse and has
larger coefficients.

The comparison is pinned to ProveIt commit
\pp{953fb4d82d4229f4efc5f9db4cc2bf38ae329bb4}.
The repository also has compact equations with far fewer variables and
much higher degree or very large fixed coefficients. At that snapshot its
complete display uses 15 strictly positive variables, 66 integer-leaf binary
operations, and degree 157 \cite{current}. That result has a different
domain and objective. Our contribution is an explicit improvement in
the direct, small-coefficient, degree-four construction.

A final read-only check found sources 160--165 at commit
\pp{fb1bed5eb7196cf35abfaa746c18c6ee7eae6cfb} \cite{late}.
Their READMEs report further tradeoffs: source 162 has a 236-variable
quartic with 4,876 monomials and height 7,968; source 165 has a
238-variable quartic with 5,236 monomials and height 3,968.
Our two forms use fewer variables, while those examples have
smaller support and height. The simultaneous improvement in
Table~\ref{tab:comparison} is specifically over the pinned
source 158 equation. Later metrics were read from the cited
READMEs, not independently expanded. Their exact file identities
are recorded in \pp{provenance/LATE_REPOSITORY_CHECK.json}.

\section{Logical calibration and the starting arithmetic system}

\subsection{The inherited chain}

The arithmetic source encodes finite derivations in the numerical calculus
\(\Ccal\), rather than a particular Turing-machine execution. Its logical
chain is
\begin{equation}\label{eq:logicalchain}
\begin{split}
 \text{integer zero}
 &\ \longleftrightarrow\ \text{finite square-shell C222 refutation}\\
 &\ \longleftrightarrow\ \text{finite original C222 refutation}\\
 &\ \longleftrightarrow\ \text{finite ZF refutation}\\
 &\ \longleftrightarrow\ \text{finite ZFC refutation}.
\end{split}
\end{equation}
The first equivalence is arithmetic; the second is a recursive recoding
of formula constructors. The last two use the declarations
\begin{quote}\small
\pp{KEndpoint.c222_iff_zf_refutation}\\
\pp{ConstructibleChoice.zf_refutation_iff_zfc_refutation}.
\end{quote}
These are unconditional calculus/refutation interfaces
\cite{endpoint,choice}; machine-specific simulation hypotheses are not used.

We retain the source statements and their pinned identities. We have
not rebuilt their entire transitive Lean dependency graph. The ordinary
proofs used here can be carried out in ZFC; this report does not supply
an encoded object-ZFC derivation or a fresh full Lean or Rocq proof term.
This is the dependency boundary of Theorem~\ref{thm:main}.

\subsection{The syntax being arithmetized}

For \(a,b\in\N\) and \(s=a+b\), the square-shell recoding uses
\[
 Q_0(a,b)=s^2+a,\qquad Q_1(a,b)=s^2+3a+4b+3
\]
and four constructors
\[
 I=2Q_0,\quad M=2Q_0+1,\quad A_{\rm syn}=2Q_1,\quad E_{\rm syn}=2Q_1+1.
\]
The two shell ranges partition the natural numbers before the parity
split. Recursive decoding yields a bijection with the original coding.
It fixes zero and the falsity code one. Constructor identities and
freshness disequalities transport through this bijection. The source
158 syntax note identifies this interface \cite{s158}.

The local compiler represents a backwards proof-tree expansion by a
finite list of polynomial alternatives. A pending multiset of labels
is represented by a product of weights \(1+jF\). Local transitions
replace a conclusion by zero, one, or two premises. A beta/CRT code
stores the accumulator path; binomial divisibility selects its rows
and enforces bounded local parameters. Pell certificates replace all
requested powers by ordinary polynomial equations.

These mechanisms are inherited. The present construction preserves the
exact local square-shell Horn circuit and changes only the identified
range and certificate layers.

\subsection{The factorial-free source}

Our actual starting system is the 242-coordinate arithmetic source in
package 157, before that package's later projections \cite{s157}.
Its literal file is a pinned input. Its semantic quantities include
\begin{equation}\label{eq:source}
\begin{gathered}
 t=16+4d^2,\qquad F=t^t,\qquad C=2c,\\
 H=2+4(C+A_1+\cdots+A_8),\qquad y=F(2v+1),\\
 B h_B=C+F+1,\\
 x=F+1+(y+1)u_0,\qquad
 x=1+[1+(t+1)y]u_t.
\end{gathered}
\end{equation}
Exponentiation in this display describes the meanings proved for
polynomial certificates; it is not a primitive of the final equation.
The initial requests are \(F=t^t\), \(N=B^x\), \(q=(BH)^x\);
the other requests use the two bases \(q\) and \(q+1\).
The binomial radius is \(F\), with authenticated digit
\(B=\binom{F+t}{t}\). Authenticating \(t!\) is unnecessary.

The retained square frame has
\begin{equation}\label{eq:frame}
 W=2+\sum_\nu a_\nu^2,\qquad W\ge258,\qquad |A|>W^W.
\end{equation}
Every actual requested base, index, and proposed output is among
the affine atoms \(a_\nu\). The retained orientation argument identifies
the Pell pairs at their canonical \emph{signed} indices before the
arithmetic signs in \eqref{eq:source} are known.
The distinction between \(a_\nu^2<W\) and the weaker \(|a_\nu|<W\)
is essential for pooling. Independent checkers reconstruct the actual
36-square identity and check all 17 retained power requests.

\section{Cofinal strict ranges}

Five old tests use lower binomial index \(t+1\): the four local
schema parameters and the derived label. We replace all five by
tests at lower index \(t\). The decoded bounds become
\begin{equation}\label{eq:strict}
 0\le a,b,c,d<t,\qquad 0\le J(a,b,c,d)<t.
\end{equation}
This is a bound on each \emph{local} tuple. We do not bound the much
larger global value of \(J\) on the CRT representatives by \(t\);
that global value remains the semantic Pell anchor.

\begin{proposition}\label{prop:strict}
The strict-range replacement preserves existence of an integer
solution of the factorial-free arithmetic source.
\end{proposition}
\begin{proof}
Soundness uses the inherited prime-power range lemma:
if \(p^e\mid\binom{a}{k}\), \(a\ge0\), \(k,e\ge1\), then the least residue of
\(a\) modulo \(p^e\) is less than \(k\). Taking \(k=t\)
strengthens the decoded local bounds, so previous soundness
conclusions continue to hold.

For completeness fix a finite refutation tree. Choose \(t=16+4d^2\)
strictly larger than its number of expansion steps, every local
parameter, and every value of the full local derived-label expression,
including inactive branches and the truth-padding tuple. This is a
cofinal choice made before constructing \(F\), the accumulators,
the beta code, or the Pell frame.

For each full prime power \(p^m\mathbin{\|}B\), hence
\(p^{xm}\mathbin{\|}N=B^x\), precision-\(2x\) CRT
makes a representative congruent to its target \(w\in[0,t-1]\)
modulo \(p^{2xm}\). If the representative is below \(t\), its
binomial at lower index \(t\) is zero. Otherwise its falling
factorial contains \(a-w\). Since
\[
 v_p(t!)\le t-1<xm,
\]
division by \(t!\) leaves divisibility by \(p^{xm}\).
Polynomial congruence gives the same conclusion for the derived label.
The factorial-free source supplies the much stronger endpoint bound
\(x\ge2+(t+1)y\), so replacing an early ratio consequence
\(x\ge t+1\) by \(x\ge t\) causes no loss.
\end{proof}

Every affected normalized digit can now use the existing \(q^t\)
denominator. Delete the \(q^{t+1}\) output, its private power
quotient, the wire for its derived Pell pair, and its normalization
ratio. Delete their three rows, redirect five ratio uses, and remove
the old output square from the frame. The system has 238 coordinates,
and all principal orientation pairs remain.

This is a semantic transformation: an old witness need not survive
at its old horizon. Completeness rebuilds witnesses at a larger
horizon. We do \emph{not} use source 158's separate
factorial-dependent change \(y=V-1\); here \(y=F(2v+1)\) is retained.

\section{Two quadratic rows for a whole block of powers}

The following fixed-base construction is inherited from source 158.
Its signed hypotheses and its use in the new bootstrap are checked
again for this combined source.

\begin{lemma}[Signed fixed-base pooling]\label{lem:pool}
Let \(W\ge258\), \(|A|>W^W\), and \(1\le r\le15\).
Let \(n\ne0\), and suppose, for each \(1\le i\le r\),
\[
 n^2<W,\qquad m_i^2<W,\qquad k_i^2<W.
\]
Let \((X_i,Y_i)\) be the canonical signed Pell pair at parameter
\(A\) and integer index \(k_i\), and put \(Z_i=AY_i\).
There exist \(K,L\in\Z\) satisfying
\begin{align}
 K-2\sum_iY_i(m_i-X_i)-(2A-n)L&=0,\label{eq:pool1}\\
 \sum_i[(m_i-X_i+Z_i)^2-Y_i^2]-nK+L&=0\label{eq:pool2}
\end{align}
if and only if every exact power identity \(m_i=n^{k_i}\) holds.
For \(|n|\ge2\), these identities force \(k_i\ge0\).
For \(n=\pm1\), signed indices are allowed.
\end{lemma}
\begin{proof}
Write
\[
 T=2An-n^2-1,\quad D_i=m_i-X_i+Z_i,\quad
 U=\sum_iY_i(m_i-X_i),\quad S=\sum_i(D_i^2-Y_i^2).
\]
The polynomial identity
\begin{equation}\label{eq:identity}
 \sum_i(D_i-nY_i)^2=S-2nU-T\sum_iY_i^2
\end{equation}
holds when \(Z_i=AY_i\). Eliminating \(K\) from the two rows gives
\(S-2nU=TL\). Thus \(T\) divides the left side of
\eqref{eq:identity}.

The Pell congruence remains valid for negative indices. In the ring
with formal square root of \(A^2-1\), sending that root to \(n-A\)
is a homomorphism modulo \(T\), because
\((n-A)^2-(A^2-1)=-T\). It sends the norm-one unit
\(A+\sqrt{A^2-1}\) to \(n\). Since \(T\equiv-1\pmod n\),
\(n\) is invertible modulo \(T\). Hence
\[
 X_i+(n-A)Y_i\equiv n^{k_i}\pmod T
\]
for all signed indices.

Put \(b=\lfloor\sqrt{W-1}\rfloor\). Multiplying by \(n^b\)
clears all negative indices and gives
\[
 T\mid E:=\sum_i(m_i n^b-n^{b+k_i})^2.
\]
All exponents are nonnegative. Every
\(|n|,|m_i|,|k_i|\le b\), with \(b\ge16\), so
\[
 0\le E\le4r b^{4b}<W^{2b+1}\le W^W.
\]
Here \(4r\le60<W\), \(b^2<W\), and
\(2b+1\le b^2+1\le W\). On the other hand,
\[
 |T|\ge2|A||n|-n^2-1\ge2|A|-W>W^W.
\]
Thus \(E=0\), so each square vanishes. If \(|n|\ge2\) and
\(k_i<0\), the resulting equation \(m_i n^{-k_i}=1\) is
impossible over \(\Z\). All other cases give the asserted powers.

Conversely the exact powers make each \(D_i-nY_i\) divisible
by \(T\). Identity \eqref{eq:identity} then makes
\(L=(S-2nU)/T\) integral. Taking \(K=2U+(2A-n)L\)
supplies both rows.
\end{proof}

The symbols \(T,U,S,b\) are proof abbreviations, not new
coordinates. The other quantities are existing affine coordinate
expressions, including the existing \(Z_i=AY_i\) gate. Thus both
new rows are quadratic. A whole block needs just two new
coordinates and two rows in this construction.

The strict system has eleven requests at \(q+1\) and three at \(q\).
Replacing individual quotients and the old modulus coordinates saves
ten and two coordinates, respectively. The initial requests
\(F=t^t\), \(N=B^x\), and \(q=(BH)^x\) remain individual.

\section{The pooled factorial-free bootstrap}\label{sec:bootstrap}

An algebraic pooling identity alone is insufficient. The bases and
signed indices must be justified in an order that avoids circular
sign assumptions.

\begin{proposition}\label{prop:bootstrap}
The modified system recovers all arithmetic signs and exact powers
without assuming a positive radix or positive CRT representatives.
\end{proposition}
\begin{proof}
The individual request at known base and index \(t\) first gives
\(F=t^t\ge16^{16}=2^{64}\). Since \(t=16+4d^2\), \(F\)
is even. With \(C=2c\), the right side of
\(Bh_B=C+F+1\) is odd even when \(c<0\). Thus \(B\) is
odd and nonzero. Independently \(H\equiv2\pmod4\), whence
\(|H|\ge2\) and \(|BH|\ge2\). The individual \(BH\) request
gives \(x\ge0\) and \(q=(BH)^x\).

The beta equation gives \(|y|\ge F\), so
\[
 |1+(t+1)y|\ge(t+1)F-1>1.
\]
If \(x=0\), the terminal endpoint makes this integer divide
\(-1\), a contradiction. Therefore \(x\ge1\) and \(|q|\ge2\).

Apply Lemma~\ref{lem:pool} to the \(q\) block. Its request at
exponent \(y+1\) forces \(y+1\ge0\). A negative \(y\) would
be at most \(-F<-1\), so \(y\ge F>0\).
The first endpoint excludes \(u_0\le-1\), giving \(x\ge F+1\);
the last endpoint now forces \(u_t\ge1\). Thus
\begin{equation}\label{eq:xlarge}
 x\ge2+(t+1)y\ge2+(t+1)F.
\end{equation}
Since \(|q|\ge2^x\), we have \(|q+1|\ge2\).
Pooling at \(q+1\) now proves all representative indices
nonnegative and all outputs exact. In particular
\(C,A_1,\ldots,A_8\ge0\), the semantic anchor retains its
meaning, and \(H>0\). Neither \(B>0\) nor \(q>0\) has
yet been assumed.

To finish the sign argument, write \(Q=|q|\) and
\(c_j=\binom{F+t}{j}\). For \(0\le j\le t\),
\[
 c_j\le(F+t)^t\le(F^2+F)^t\le2^{t(F-1)}.
\]
Equation \eqref{eq:xlarge} yields
\(c_{\max}/(Q-1)<1/32\) and \(c_t<Q/32\).
The binomial theorem is valid at negative \(q\), and its true low
part \(r_0=\sum_{j<t}c_jq^j\) satisfies \(|r_0|<Q^t/32\).

The retained normalization identities and common remainder norm
bound each supplied remainder by half its denominator in absolute
value. Congruent supplied and genuine remainders must therefore
coincide. The digit row gives \(B\equiv c_t\pmod q\).
For \(a=|BH|\ge2\), \(|B|\le a<Q/2\), because \(x\ge3\).
Hence \(|B-c_t|<Q\), and
\[
 B=c_t=\binom{F+t}{t}>0.
\]
Now \(q>0\), \(N=B^x>0\), and all later binomial, CRT, and
weight arguments have their original natural-number meanings.
\end{proof}

Source 157's remaining estimates are unchanged. They use
\(\operatorname{lcm}(1,\ldots,t)^2\le16^{t-1}<t^t=F\)
for binomial CRT selection, its sharper weight-separation proof
at \(F=t^t\), and its local polynomial majorant below \(2^x\)
\cite{s157}. The local ranges have only become smaller.

For reverse construction choose the larger cofinal horizon,
choose the beta spacing as a sufficiently large odd multiple of
the odd part of \(t!\) times \(F\), then build the original CRT
and simultaneous Pell witnesses. The converse of
Lemma~\ref{lem:pool} supplies the four pooled coordinates.
Old witnesses need not remain numerically unchanged.

\section{Exact algebraic compression}

\subsection{Private wires and Pell difference coordinates}

Before pooling, three private coordinates are removed by the
consumer-based integer changes of source 158. Their forward and
reverse residual identities are checked against the factorial-free
source by matching the unchanged local Horn subcircuit.
This takes 238 coordinates to 235. The change introduces one
shared product, eliminates four private wires, and replaces a gate
by an aggregate square-sum coordinate. It is an exact algebraic
transformation, unlike the cofinal change of Section~3.

Pooling gives 223 coordinates and 160 residuals.
For each retained Pell pair use the difference coordinate
\[
 U_i=X_i-Z_i,\qquad X_i=U_i+Z_i.
\]
Because \(Z_i\) already has a gate coordinate, this is an invertible
integer affine change in the ambient coordinate system.
On the common zero set, the retained gate equation \(Z_i=AY_i\)
identifies \(U_i\) with \(X_i-AY_i\).
It reduces expanded support without changing the number of
unknowns or requiring a divisibility premise. The resulting
223-coordinate system is the pinned input to the following
projection theorem.

\subsection{Removing common quadratic directions}

\begin{theorem}[Integral radical projection]\label{thm:radical}
Let \(F_0,\ldots,F_{m-1}\in\Z[\mathbf{x}]\) have degree at most two.
Suppose \(v\in\Z^n\) satisfies
\[
 F_i(\mathbf{x}+uv)=F_i(\mathbf{x})+u a_i
 \quad\text{for every }i,
\]
where \(a_i\in\Z\) is independent of \(\mathbf{x}\).
Assume \(v_p=1\) and \(a_r=\pm1\).
On \(z_p=0\), set
\[
 G_i(\mathbf{z})=F_i(\mathbf{z})-(a_i/a_r)F_r(\mathbf{z})
 \qquad(i\ne r).
\]
The common integer zero sets of \(F\) and \(G\) are in polynomial
bijection, and every \(G_i\) has degree at most two.
\end{theorem}
\begin{proof}
The projection and lift are
\[
 \pi(\mathbf{x})=\mathbf{x}-x_pv,\qquad
 \lambda(\mathbf{z})=\mathbf{z}-vF_r(\mathbf{z})/a_r.
\]
They have integer coefficients because \(a_r=\pm1\).
The directional identity gives
\[
 F_r(\lambda(\mathbf{z}))=0,\qquad
 F_i(\lambda(\mathbf{z}))=G_i(\mathbf{z}).
\]
At a zero of \(F\),
\(F_r(\pi(\mathbf{x}))=-a_rx_p\), and therefore
\(\lambda(\pi(\mathbf{x}))=\mathbf{x}\).
Conversely \(z_p=0\) gives
\(\pi(\lambda(\mathbf{z}))=\mathbf{z}\).
\end{proof}

A common kernel vector for the homogeneous Hessian matrices has
the required constant-direction property. The certificate supplies
five integer vectors, pivots, omitted rows, and directional
constants. The checker verifies the directional identities
themselves, not merely a reported rank.

\begin{table}[ht]
\centering\small
\begin{tabular}{rrrl}
\toprule
Dimension& Pivot & Omitted row & Nonzero entries of \(v\)\\
\midrule
223&43&28&\(v_{43}=1,\ v_{44}=-1\)\\
222&59&45&\(v_{59}=1,\ v_{60}=-1,\ v_{61}=-2\)\\
221&63&48&\(v_{61}=-1,\ v_{63}=1\)\\
220&30&15&\(v_{30}=1,\ v_{31}=-1,\ v_{32}=-2\)\\
219&24&9 &\(v_{24}=1,\ v_{26}=-1,\ v_{27}=-2\)\\
\bottomrule
\end{tabular}
\caption{The five actual radical projections. Indices refer to the
current system before that row of the table is applied.}
\label{tab:directions}
\end{table}

The transformations take 223 coordinates and 160 rows to 218
coordinates and 155 rows. This is source 159's method applied to
the new combined arithmetic source \cite{s159}; the generic
projection theorem is not claimed as a new literature theorem.

Then 38 recorded integer row operations lower height and support.
Each \(F_i\leftarrow F_i+kF_j\), \(i\ne j\), has integer inverse
\(F_i\leftarrow F_i-kF_j\), so the common zero set is unchanged.
The sum of squares of the final rows is \(P_{218}\).

\subsection{The 216-coordinate companion}

Two affine consequences permit two additional coordinate
eliminations. Each certificate expresses an affine polynomial as
an exact constant-coefficient combination of the rows, with unit
coefficient on the removed coordinate. Solving for that coordinate
and substituting preserves degree two. The certificate identifies
the omitted old row and proves its recovery from the combination.
The row-combination coefficient need only be nonzero for this
recovery over \(\Z\); the solved coordinate coefficient must be a
unit to ensure an integral lift.

Direct substitution needlessly enlarges coefficients. In the
218-coordinate projected system before the balancing row operations,
let
\[
 a=x_8,\quad b=x_9,\quad u=x_{19},\quad
 v=x_{20},\quad w=x_{33}.
\]
The exact identity is
\[
 -F_5+F_{18}=4+12a+16b+4u-v+w.
\]
Introduce the integral affine coordinates
\[
 t_*=u+3a+4b+1,\qquad c_*=v-2t_*,
\]
whose inverse is
\[
 u=t_*-3a-4b-1,\qquad v=c_*+2t_*.
\]
The affine relation becomes \(2t_*-c_*+w=0\), hence
\(w=c_*-2t_*\). The retained common-square row collapses to
\[
 c_*-4t_*^2+1=0.
\]
No integer coordinate is divided by two. The actual chart and
inverse are explicit certificate data, and the notation here does
not add polynomial unknowns.

The two affine eliminations, followed by 59 invertible row
operations, yield \(P_{216}\). Its 153 quadratic residuals and full
collected expansion are included. The centered chart controls the
coefficient growth caused by the two additional substitutions.
The principal \(P_{218}\) remains the better sparse, low-height
choice.

\section{Proof of the main theorem}

\begin{proof}[Proof of Theorem~\ref{thm:main}]
A finite sum of integer squares is zero precisely when every
residual is zero. Thus a zero of \(P_{218}\) is a zero of its
155-row quadratic system.

Invert the recorded row operations and apply the lifts of
Theorem~\ref{thm:radical}, obtaining a 223-coordinate zero.
Invert the Pell difference chart. The actual square-frame
memberships, Lemma~\ref{lem:pool}, and the bootstrap of
Proposition~\ref{prop:bootstrap} restore every exact power.
The exact private-wire maps and the strict-range soundness
argument recover the inherited local proof-certificate
interpretation. It gives a finite square-shell C222 refutation.
The recoding and named interfaces in \eqref{eq:logicalchain}
produce a finite ZFC refutation.

Conversely a finite ZFC refutation yields a finite C222
refutation and its square-shell recoding by the inherited
calibration. Choose a sufficiently large cofinal horizon and
construct the factorial-free strict-range witnesses, as in
Section~3 and source 157's completeness proof.
The private-wire changes, the converse of
Lemma~\ref{lem:pool}, and the Pell difference chart give a
223-coordinate zero. Apply the radical projections and
recorded row operations to obtain a zero of \(P_{218}\).

All steps are finite integer arithmetic, recursive syntactic
constructions, or the identified inherited calibration, and
their arguments formalize in ZFC. Negation gives the stated
consistency equivalence. The additional centered chart and
affine projections give the same equivalence for \(P_{216}\).
\end{proof}

\section{Finite verification and reproducibility}

Verification separates three claims. The article and independent
mathematical review establish the new unbounded semantic steps
using the identified inherited arithmetic theorems.
The composition checkers bind those statements to the literal
source edits, all actual power requests and the square frame,
and the private coordinate maps. The algebra checker verifies
every projection, row operation, inverse chart, and collected
coefficient. Coefficient checks do not replace unbounded proofs.

The root commands use only Python's standard library:
\begin{verbatim}
python3 -B verify.py
python3 -B -O verify.py
\end{verbatim}
Both ordinary and optimized execution are supported; required
checks are not removable Python assertions. The receipts record
the checks actually performed. Deliberately corrupted
certificates are rejected. There is no search for a zero.

The sparse format stores each term as an integer coefficient
and a sorted list of coordinate indices. Repetitions mean
repeated factors; the empty list denotes the constant monomial.
For example \([{-}3,[2,7,7]]\) denotes \(-3x_2x_7^2\).
Monomial index lists in residuals have length at most two, and
in the expanded quartic at most four. Validation excludes duplicate monomials, zero
coefficients, noninteger data, and out-of-range indices.

The verifier independently expands \(\sum R_j^2\), compares
every coefficient with the collected polynomial, and checks
degree, support, height, coordinate liveness, and the printed
principal appendix. All 155 residuals in that appendix are
parsed back into integer polynomials for comparison.
Raw source identities and canonical JSON hashes are recorded
separately.

The complete literal equations can be evaluated without
executing a historical compiler. The builder and frozen
certificates are also supplied for reproduction.
The following table gives the construction stages; changes
of integer basis leave the counts of variables and rows fixed.

\begin{center}\small
\begin{tabular}{lrr}
\toprule
Stage&Coordinates&Quadratic rows\\
\midrule
Pinned factorial-free arithmetic source&242&178\\
Strict local ranges&238&175\\
Private-wire basis&235&172\\
Pooled powers and Pell differences&223&160\\
Common-direction projection&218&155\\
Centered affine projections&216&153\\
\bottomrule
\end{tabular}
\end{center}

\subsection{Attribution and verification boundary}

The new result is the complete cross-branch construction, its
checked noncircular bootstrap, the exact finite bindings, and
the two final literal polynomials. The factorial-free
binomial/CRT theorem is inherited from source 157; the pooling
lemma and private-wire identities from source 158; the
common-direction projection method from source 159.
The combined instance and its improved equations are the
contribution of this iteration.

New finite checks were executed for this release. Historical
source receipts remain historical. The full logical formal
development was not recompiled. No global optimum or
literature-wide priority is asserted, and the internal
independent reviews are AI-assisted audits, not external
human peer review.

\section{Further research questions}

\begin{enumerate}
\item Can a residual basis for the centered 216-coordinate
companion approach the support and height of \(P_{218}\)?
This is now a finite integer-algebra problem on a fixed
coefficient system.
\item Can a pooled coordinate be eliminated while all rows
stay quadratic and the signed bootstrap remains valid?
The two-row lemma is a construction, not a minimum theorem.
\item Can the local square-shell proof circuit be redesigned
to retain more common quadratic directions through compilation,
while preserving the constructor recoding and freshness guards?
\item Can low-degree consequences beyond constant row
combinations and common Hessian kernels expose further
integral projections? Rational eliminations need integral
lifts or separate divisibility arguments.
\item A focused formalization could first establish
Lemma~\ref{lem:pool} and Theorem~\ref{thm:radical}, and then
validate a compact certificate format. That isolates the new
mathematics without requiring a new full logical build.
\end{enumerate}

These questions are not assumptions of either delivered equation.

\appendix
\section{Every residual of the principal equation}
\label{app:residuals}

The following list fixes all coefficients of
\(P_{218}=\sum_{j=0}^{154}R_j^2\).
The only independent unknowns are
\(x_0,\ldots,x_{217}\in\Z\).
There are no omitted constants, variable exponents, range
predicates, or additional side conditions.

\begingroup\footnotesize
\setlength{\jot}{1.5pt}
% Literal residual appendix. Generated from algebra/equation218.json.
\begin{align*}
R_{0} &= x_{0} -x_{0}\,x_{2} -x_{2} +x_{7}
\end{align*}
\begin{align*}
R_{1} &= -x_{2}\,x_{5} -x_{3} +x_{6}\,x_{7}
\end{align*}
\begin{align*}
R_{2} &= -1 +x_{1} -x_{6}\,x_{7}
\end{align*}
\begin{align*}
R_{3} &= 1 +x_{3} -x_{4}\,x_{18} +2x_{16}
\end{align*}
\begin{align*}
R_{4} &= -1 +2x_{8} -x_{8}^{2} -2x_{8}\,x_{9} +2x_{9} -x_{9}^{2} +x_{19}
\end{align*}
\begin{align*}
R_{5} &= -4 -12x_{8} -16x_{9} -4x_{19} +x_{20} -x_{33}
\end{align*}
\begin{align*}
R_{6} &= -2 +2x_{8} -x_{8}^{2} -4x_{8}\,x_{20} +4x_{20} -4x_{20}^{2} +x_{21}
\end{align*}
\begin{align*}
R_{7} &= 67 -60x_{8} -3x_{8}^{2} -60x_{8}\,x_{20} -24x_{8}\,x_{21} -116x_{20} -124x_{20}^{2} -64x_{20}\,x_{21} \\
&\quad {} -31x_{21} -4x_{21}^{2} +x_{22}
\end{align*}
\begin{align*}
R_{8} &= -2 -x_{9}^{2} -2x_{9}\,x_{10} -x_{10}^{2} +x_{23}
\end{align*}
\begin{align*}
R_{9} &= -1 +2x_{8} -x_{8}^{2} -2x_{8}\,x_{10} +2x_{10} -x_{10}^{2} +x_{24}
\end{align*}
\begin{align*}
R_{10} &= 15 -38x_{8} -29x_{8}^{2} -34x_{8}\,x_{9} -36x_{8}\,x_{10} -24x_{8}\,x_{19} -24x_{8}\,x_{24} +25x_{9} +38x_{9}^{2} \\
&\quad {} -48x_{9}\,x_{10} +16x_{9}\,x_{23} -28x_{9}\,x_{24} +36x_{10} +70x_{10}^{2} -32x_{10}\,x_{19} +32x_{10}\,x_{23} \\
&\quad {} -17x_{19} -8x_{19}\,x_{24} +40x_{23} -14x_{24} +x_{25} -x_{29} -x_{33}
\end{align*}
\begin{align*}
R_{11} &= -15 +14x_{8} -2x_{8}^{2} -8x_{8}\,x_{9} -8x_{8}\,x_{23} +4x_{9}^{2} +8x_{9}\,x_{23} +4x_{23}^{2} -4x_{25}^{2} \\
&\quad {} +x_{26} -3x_{44}
\end{align*}
\begin{align*}
R_{12} &= -6 +16x_{10} -16x_{10}^{2} -16x_{10}\,x_{23} -16x_{10}\,x_{24} +8x_{23} -4x_{23}^{2} -8x_{23}\,x_{24} +8x_{24} \\
&\quad {} -4x_{24}^{2} +x_{27}
\end{align*}
\begin{align*}
R_{13} &= -6 -48x_{8}\,x_{10} -48x_{8}\,x_{24} -24x_{8}\,x_{27} -64x_{9}\,x_{10} -64x_{9}\,x_{24} -32x_{9}\,x_{27} +7x_{10} \\
&\quad {} +9x_{10}^{2} -16x_{10}\,x_{19} +8x_{10}\,x_{24} +4x_{10}\,x_{27} -16x_{19}\,x_{24} -8x_{19}\,x_{27} +2x_{20} +8x_{24} \\
&\quad {} +4x_{27} +x_{28} -3x_{33} -x_{43}
\end{align*}
\begin{align*}
R_{14} &= 6 -19x_{8} -35x_{8}^{2} -80x_{8}\,x_{10} +4x_{8}\,x_{23} -24x_{8}\,x_{24} +33x_{9} +35x_{9}^{2} +80x_{9}\,x_{10} \\
&\quad {} +24x_{9}\,x_{23} -4x_{9}\,x_{24} +16x_{10} +32x_{10}\,x_{23} -32x_{10}\,x_{24} +12x_{23} +4x_{23}^{2} -8x_{24} \\
&\quad {} -4x_{24}^{2} +x_{29} +4x_{29}^{2} -x_{30}
\end{align*}
\begin{align*}
R_{15} &= 90 -10x_{8} -2x_{8}^{2} -24x_{8}\,x_{9} -32x_{8}\,x_{10} -8x_{8}\,x_{23} -72x_{9} -18x_{9}^{2} -84x_{9}\,x_{10} \\
&\quad {} -48x_{9}\,x_{23} -96x_{10} -74x_{10}^{2} -64x_{10}\,x_{23} -78x_{23} -8x_{23}^{2} -4x_{29}^{2} +x_{30}
\end{align*}
\begin{align*}
R_{16} &= -12 +16x_{8} -4x_{8}^{2} -8x_{8}\,x_{10} -8x_{8}\,x_{19} -8x_{8}\,x_{23} +16x_{10} -4x_{10}^{2} -8x_{10}\,x_{19} \\
&\quad {} -8x_{10}\,x_{23} +16x_{19} -4x_{19}^{2} -8x_{19}\,x_{23} +16x_{23} -4x_{23}^{2} +x_{31}
\end{align*}
\begin{align*}
R_{17} &= 25 -56x_{8} -48x_{8}^{2} -128x_{8}\,x_{10} -64x_{8}\,x_{23} -32x_{8}\,x_{24} -32x_{8}\,x_{31} -90x_{10} \\
&\quad {} -52x_{10}^{2} -48x_{10}\,x_{23} -8x_{10}\,x_{24} -40x_{10}\,x_{31} -24x_{23} -8x_{23}^{2} -16x_{23}\,x_{31} \\
&\quad {} -38x_{24} +4x_{24}^{2} -8x_{24}\,x_{31} -2x_{27} -4x_{31} -4x_{31}^{2} +x_{32}
\end{align*}
\begin{align*}
R_{18} &= 8 +6x_{8} -8x_{8}^{2} -64x_{8}\,x_{9} +16x_{8}\,x_{19} -64x_{9}^{2} +6x_{19} +8x_{19}^{2} -x_{20} +3x_{33} \\
&\quad {} -x_{37}
\end{align*}
\begin{align*}
R_{19} &= -2 +2x_{8} -x_{8}^{2} -4x_{8}\,x_{33} +4x_{33} -4x_{33}^{2} +x_{34}
\end{align*}
\begin{align*}
R_{20} &= 67 -60x_{8} -3x_{8}^{2} -60x_{8}\,x_{33} -24x_{8}\,x_{34} -116x_{33} -124x_{33}^{2} -64x_{33}\,x_{34} \\
&\quad {} -31x_{34} -4x_{34}^{2} +x_{35}
\end{align*}
\begin{align*}
R_{21} &= -x_{9}^{2} -4x_{9}\,x_{35} -4x_{35}^{2} +x_{36}
\end{align*}
\begin{align*}
R_{22} &= -1 -28x_{8}^{2} -32x_{8}\,x_{9} -40x_{8}\,x_{19} +8x_{9} -32x_{9}\,x_{19} -4x_{19} -12x_{19}^{2} -x_{20} +x_{37}
\end{align*}
\begin{align*}
R_{23} &= 49 -34x_{9}^{2} -44x_{9}\,x_{10} -24x_{9}\,x_{23} -32x_{9}\,x_{37} +6x_{10} -10x_{10}^{2} -24x_{10}\,x_{23} \\
&\quad {} -24x_{10}\,x_{37} -32x_{23} -8x_{23}\,x_{37} +4x_{37} -4x_{37}^{2} +x_{38} -x_{44}
\end{align*}
\begin{align*}
R_{24} &= 8x_{8} -64x_{8}^{2} -96x_{8}\,x_{10} -32x_{8}\,x_{24} -32x_{8}\,x_{38} +6x_{10} -36x_{10}^{2} -24x_{10}\,x_{24} \\
&\quad {} -24x_{10}\,x_{38} +2x_{24} -4x_{24}^{2} -8x_{24}\,x_{38} +4x_{38} -4x_{38}^{2} +x_{39}
\end{align*}
\begin{align*}
R_{25} &= -x_{10}^{2} -2x_{10}\,x_{11} -x_{11}^{2} +x_{40}
\end{align*}
\begin{align*}
R_{26} &= -11 +3x_{8}^{2} +8x_{8}\,x_{10} -8x_{8}\,x_{24} -4x_{8}\,x_{40} +4x_{9}^{2} +8x_{9}\,x_{23} -8x_{9}\,x_{24} -8x_{10} \\
&\quad {} +2x_{10}^{2} -8x_{10}\,x_{40} +4x_{23}^{2} -8x_{23}\,x_{24} +18x_{24} -4x_{24}^{2} +4x_{25}^{2} -x_{26} +4x_{40} \\
&\quad {} -4x_{40}^{2} +x_{41}
\end{align*}
\begin{align*}
R_{27} &= 4x_{8} -18x_{8}^{2} -40x_{8}\,x_{10} -32x_{8}\,x_{24} -40x_{8}\,x_{40} -16x_{8}\,x_{41} +4x_{10} -24x_{10}^{2} \\
&\quad {} -32x_{10}\,x_{24} -48x_{10}\,x_{40} -16x_{10}\,x_{41} -16x_{24}^{2} -32x_{24}\,x_{40} -16x_{24}\,x_{41} +4x_{40} \\
&\quad {} -24x_{40}^{2} -16x_{40}\,x_{41} -4x_{41}^{2} +x_{42}
\end{align*}
\begin{align*}
R_{28} &= x_{10} -25x_{10}^{2} -40x_{10}\,x_{24} -20x_{10}\,x_{27} -16x_{24}^{2} -16x_{24}\,x_{27} -4x_{27}^{2} +x_{43}
\end{align*}
\begin{align*}
R_{29} &= -1 +8x_{9} -4x_{9}^{2} -8x_{9}\,x_{23} +8x_{23} -4x_{23}^{2} +x_{44}
\end{align*}
\begin{align*}
R_{30} &= -1 -16x_{9}^{2} -8x_{9}\,x_{10} -32x_{9}\,x_{23} -16x_{9}\,x_{44} -x_{10}^{2} -8x_{10}\,x_{23} -4x_{10}\,x_{44} \\
&\quad {} -16x_{23}^{2} -16x_{23}\,x_{44} -4x_{44}^{2} +x_{45}
\end{align*}
\begin{align*}
R_{31} &= 3 +x_{9} -2x_{9}\,x_{10} -4x_{9}\,x_{45} -3x_{10}^{2} -8x_{10}\,x_{45} -x_{23} -4x_{45}^{2} +x_{46}
\end{align*}
\begin{align*}
R_{32} &= -1 -4x_{46} -4x_{46}^{2} +x_{47}
\end{align*}
\begin{align*}
R_{33} &= -1 -x_{10}^{2} -8x_{10}\,x_{46} -4x_{10}\,x_{47} -16x_{46}^{2} -16x_{46}\,x_{47} -4x_{47}^{2} +x_{48}
\end{align*}
\begin{align*}
R_{34} &= 4 +x_{9} -2x_{9}\,x_{10} -4x_{9}\,x_{48} -3x_{10}^{2} -8x_{10}\,x_{48} -x_{23} -4x_{48}^{2} +x_{49}
\end{align*}
\begin{align*}
R_{35} &= -4x_{43}^{2} -8x_{43}\,x_{49} -4x_{49}^{2} +x_{50}
\end{align*}
\begin{align*}
R_{36} &= -9 +7x_{8} -x_{8}^{2} -8x_{8}\,x_{43} -4x_{8}\,x_{50} +24x_{43} -16x_{43}^{2} -16x_{43}\,x_{50} +12x_{50} \\
&\quad {} -4x_{50}^{2} +x_{51}
\end{align*}
\begin{align*}
R_{37} &= -1 -4x_{51} -4x_{51}^{2} +x_{52}
\end{align*}
\begin{align*}
R_{38} &= -x_{10}^{2} -8x_{10}\,x_{51} -4x_{10}\,x_{52} -16x_{51}^{2} -16x_{51}\,x_{52} -4x_{52}^{2} +x_{53}
\end{align*}
\begin{align*}
R_{39} &= -x_{8} -x_{8}^{2} -4x_{8}\,x_{10} -4x_{8}\,x_{53} -4x_{10} -4x_{10}^{2} -8x_{10}\,x_{53} -4x_{53} -4x_{53}^{2} +x_{54}
\end{align*}
\begin{align*}
R_{40} &= -10 -12x_{54} -4x_{54}^{2} +x_{55}
\end{align*}
\begin{align*}
R_{41} &= 2 -4x_{9} -16x_{9}\,x_{54} -8x_{9}\,x_{55} -4x_{23} -16x_{23}\,x_{54} -8x_{23}\,x_{55} -x_{44} +8x_{54} \\
&\quad {} -16x_{54}^{2} -16x_{54}\,x_{55} +4x_{55} -4x_{55}^{2} +x_{56}
\end{align*}
\begin{align*}
R_{42} &= -10 +x_{9} -25x_{9}^{2} -40x_{9}\,x_{23} -20x_{9}\,x_{56} -16x_{23}^{2} -16x_{23}\,x_{56} -12x_{49} -4x_{49}^{2} \\
&\quad {} -4x_{56}^{2} +x_{57}
\end{align*}
\begin{align*}
R_{43} &= 21 +24x_{49} -8x_{49}^{2} -16x_{49}\,x_{57} -2x_{57} -4x_{57}^{2} +x_{58}
\end{align*}
\begin{align*}
R_{44} &= 1 +4x_{8} -4x_{8}^{2} -8x_{8}\,x_{19} -8x_{8}\,x_{60} +x_{10} -x_{10}^{2} -4x_{10}\,x_{58} +4x_{19} -4x_{19}^{2} \\
&\quad {} -8x_{19}\,x_{60} -4x_{58}^{2} +x_{59} +4x_{60} -4x_{60}^{2}
\end{align*}
\begin{align*}
R_{45} &= -4x_{8} -4x_{8}^{2} -4x_{8}\,x_{10} -8x_{8}\,x_{19} -x_{10} -x_{10}^{2} -4x_{10}\,x_{19} -4x_{19} -4x_{19}^{2} +x_{60}
\end{align*}
\begin{align*}
R_{46} &= 4x_{8} -24x_{8}^{2} -48x_{8}\,x_{19} -16x_{8}\,x_{59} -16x_{8}\,x_{60} +4x_{19} -24x_{19}^{2} -16x_{19}\,x_{59} \\
&\quad {} -16x_{19}\,x_{60} -4x_{59}^{2} +8x_{60} -8x_{60}^{2} +x_{61}
\end{align*}
\begin{align*}
R_{47} &= -2 -4x_{61} -4x_{61}^{2} +x_{62}
\end{align*}
\begin{align*}
R_{48} &= -x_{3}\,x_{13} +x_{63}
\end{align*}
\begin{align*}
R_{49} &= x_{12} -x_{13} -2x_{22}\,x_{63} -x_{64}
\end{align*}
\begin{align*}
R_{50} &= x_{12} -x_{13} -2x_{26}\,x_{63} -x_{65}
\end{align*}
\begin{align*}
R_{51} &= -2x_{28}\,x_{63} +2x_{42}\,x_{63} -2x_{63} -x_{66} +x_{71}
\end{align*}
\begin{align*}
R_{52} &= x_{12} -x_{13} -2x_{30}\,x_{63} -x_{67}
\end{align*}
\begin{align*}
R_{53} &= x_{12} -x_{13} -2x_{32}\,x_{63} +6x_{63} -x_{68}
\end{align*}
\begin{align*}
R_{54} &= x_{12} -x_{13} -4x_{35}\,x_{63} -2x_{36}\,x_{63} -x_{69}
\end{align*}
\begin{align*}
R_{55} &= x_{12} -x_{13} -2x_{39}\,x_{63} -x_{70}
\end{align*}
\begin{align*}
R_{56} &= x_{12} -x_{13} -2x_{42}\,x_{63} -4x_{63} -x_{71}
\end{align*}
\begin{align*}
R_{57} &= -x_{8}\,x_{9} +x_{8}\,x_{10} +x_{9} +x_{9}\,x_{10} -x_{10} -x_{10}^{2} +x_{72}
\end{align*}
\begin{align*}
R_{58} &= 1 +2x_{11} +x_{11}^{2} -x_{72}^{2} +x_{73}
\end{align*}
\begin{align*}
R_{59} &= -x_{12}^{2} +2x_{12}\,x_{13} +2x_{12}\,x_{75} -x_{13}^{2} -2x_{13}\,x_{75} -x_{73}^{2} +x_{74} -x_{75}^{2}
\end{align*}
\begin{align*}
R_{60} &= -4x_{61}\,x_{63} -2x_{62}\,x_{63} +2x_{63} +x_{75}
\end{align*}
\begin{align*}
R_{61} &= -x_{9}\,x_{63} -x_{13} +x_{76}\,x_{77} +x_{78}
\end{align*}
\begin{align*}
R_{62} &= -1 +2x_{3} -2x_{3}\,x_{8} -2x_{3}\,x_{19} +x_{76}
\end{align*}
\begin{align*}
R_{63} &= -x_{8}\,x_{213} -x_{12} +x_{77} +x_{213}
\end{align*}
\begin{align*}
R_{64} &= -6x_{8}\,x_{63} -8x_{9}\,x_{63} +x_{9}\,x_{213} +x_{12} -x_{13} -2x_{19}\,x_{63} +x_{79}
\end{align*}
\begin{align*}
R_{65} &= -x_{12}\,x_{64} +x_{13}\,x_{64} +x_{80}
\end{align*}
\begin{align*}
R_{66} &= -x_{65}\,x_{80} +x_{81}
\end{align*}
\begin{align*}
R_{67} &= -x_{66}\,x_{81} +x_{82}
\end{align*}
\begin{align*}
R_{68} &= -x_{67}\,x_{82} +x_{83}
\end{align*}
\begin{align*}
R_{69} &= -x_{68}\,x_{83} +x_{84}
\end{align*}
\begin{align*}
R_{70} &= -x_{69}\,x_{84} +x_{85}
\end{align*}
\begin{align*}
R_{71} &= -x_{70}\,x_{85} +x_{86}
\end{align*}
\begin{align*}
R_{72} &= -x_{71}\,x_{86} +x_{87}
\end{align*}
\begin{align*}
R_{73} &= -x_{78}\,x_{87} +x_{88}
\end{align*}
\begin{align*}
R_{74} &= -x_{79}\,x_{88} +x_{89}
\end{align*}
\begin{align*}
R_{75} &= -x_{89}^{2} +x_{90}
\end{align*}
\begin{align*}
R_{76} &= -x_{2} -2x_{2}\,x_{16} +2x_{16} +x_{91}
\end{align*}
\begin{align*}
R_{77} &= -x_{14}\,x_{91} +x_{92}
\end{align*}
\begin{align*}
R_{78} &= -2x_{1}^{2} +2x_{1}\,x_{12} +2x_{1}\,x_{13} +2x_{1}\,x_{92} +2x_{1}\,x_{93} -x_{12}^{2} -2x_{12}\,x_{92} -x_{13}^{2} \\
&\quad {} -2x_{13}\,x_{93} +x_{17}\,x_{95} -x_{74}\,x_{90} -x_{92}^{2} -x_{93}^{2}
\end{align*}
\begin{align*}
R_{79} &= -x_{2}\,x_{15} +x_{15} -x_{15}\,x_{91} +x_{93}
\end{align*}
\begin{align*}
R_{80} &= -9x_{8} -9x_{9} -4x_{19} -2x_{22} -2x_{26} -2x_{28} -2x_{30} -2x_{32} -4x_{35} -2x_{36} -2x_{39} \\
&\quad {} -2x_{42} -4x_{61} -2x_{62} +x_{94}
\end{align*}
\begin{align*}
R_{81} &= -x_{96}\,x_{99} +x_{99} +x_{101}
\end{align*}
\begin{align*}
R_{82} &= -x_{4}\,x_{122} +x_{97}\,x_{204} -x_{100} -x_{101}\,x_{122}
\end{align*}
\begin{align*}
R_{83} &= x_{102}\,x_{204} -x_{103}\,x_{125} -x_{104}
\end{align*}
\begin{align*}
R_{84} &= x_{105}\,x_{204} -x_{106}\,x_{125} -x_{107}
\end{align*}
\begin{align*}
R_{85} &= x_{108}\,x_{204} -x_{109}\,x_{125} -x_{110}
\end{align*}
\begin{align*}
R_{86} &= x_{111}\,x_{204} -x_{112}\,x_{125} -x_{113}
\end{align*}
\begin{align*}
R_{87} &= x_{114}\,x_{205} -x_{116}\,x_{125} -x_{117}
\end{align*}
\begin{align*}
R_{88} &= x_{118}\,x_{205} -x_{119}\,x_{125} -x_{120}
\end{align*}
\begin{align*}
R_{89} &= -x_{17}\,x_{122} +x_{125}
\end{align*}
\begin{align*}
R_{90} &= x_{121} -x_{123}\,x_{125} -x_{124}
\end{align*}
\begin{align*}
R_{91} &= -x_{125}\,x_{127} +x_{126} -x_{128}
\end{align*}
\begin{align*}
R_{92} &= -x_{125}\,x_{130} +x_{129}\,x_{204} -x_{131}
\end{align*}
\begin{align*}
R_{93} &= -x_{125}\,x_{133} +x_{132}\,x_{204} -x_{134}
\end{align*}
\begin{align*}
R_{94} &= x_{135} -2x_{135}\,x_{137} +x_{138}
\end{align*}
\begin{align*}
R_{95} &= -x_{135}\,x_{138} -x_{138} +x_{139}
\end{align*}
\begin{align*}
R_{96} &= 2x_{136} +x_{136}^{2} +x_{138}^{2} -x_{139}^{2}
\end{align*}
\begin{align*}
R_{97} &= -x_{136}\,x_{141} -x_{141} +x_{142}
\end{align*}
\begin{align*}
R_{98} &= -1 +x_{140}^{2} +2x_{140}\,x_{142} +x_{141}^{2}
\end{align*}
\begin{align*}
R_{99} &= 1 -2x_{0} +x_{0}^{2} -2x_{0}\,x_{136} +x_{143}
\end{align*}
\begin{align*}
R_{100} &= -x_{0}\,x_{141} +x_{3} -x_{140} +x_{143}\,x_{144}
\end{align*}
\begin{align*}
R_{101} &= 1 -2x_{4} +x_{4}^{2} -2x_{4}\,x_{136} +x_{145}
\end{align*}
\begin{align*}
R_{102} &= -x_{4}\,x_{177} +x_{17} +x_{145}\,x_{146} -x_{176}
\end{align*}
\begin{align*}
R_{103} &= -1 +2x_{136}\,x_{188} -x_{147} +2x_{188} -x_{188}^{2}
\end{align*}
\begin{align*}
R_{104} &= -1 +x_{96} +x_{147}\,x_{148} -x_{176} -x_{177}\,x_{188}
\end{align*}
\begin{align*}
R_{105} &= -x_{136}\,x_{150} -x_{150} +x_{151}
\end{align*}
\begin{align*}
R_{106} &= -1 +x_{149}^{2} +2x_{149}\,x_{151} +x_{150}^{2}
\end{align*}
\begin{align*}
R_{107} &= 1 +x_{136} -x_{136}\,x_{153} -x_{153} +x_{154}
\end{align*}
\begin{align*}
R_{108} &= x_{152}^{2} +2x_{152}\,x_{154} -2x_{153} +x_{153}^{2}
\end{align*}
\begin{align*}
R_{109} &= -x_{136}\,x_{156} -x_{156} +x_{157}
\end{align*}
\begin{align*}
R_{110} &= -1 +x_{155}^{2} +2x_{155}\,x_{157} +x_{156}^{2}
\end{align*}
\begin{align*}
R_{111} &= -x_{136}\,x_{159} -x_{159} +x_{160}
\end{align*}
\begin{align*}
R_{112} &= -1 +x_{158}^{2} +2x_{158}\,x_{160} +x_{159}^{2}
\end{align*}
\begin{align*}
R_{113} &= -x_{136}\,x_{162} -x_{162} +x_{163}
\end{align*}
\begin{align*}
R_{114} &= -1 +x_{161}^{2} +2x_{161}\,x_{163} +x_{162}^{2}
\end{align*}
\begin{align*}
R_{115} &= -x_{136}\,x_{165} -x_{165} +x_{166}
\end{align*}
\begin{align*}
R_{116} &= -1 +x_{164}^{2} +2x_{164}\,x_{166} +x_{165}^{2}
\end{align*}
\begin{align*}
R_{117} &= -x_{136}\,x_{168} -x_{168} +x_{169}
\end{align*}
\begin{align*}
R_{118} &= -1 +x_{167}^{2} +2x_{167}\,x_{169} +x_{168}^{2}
\end{align*}
\begin{align*}
R_{119} &= -x_{136}\,x_{171} -x_{171} +x_{172}
\end{align*}
\begin{align*}
R_{120} &= -1 +x_{170}^{2} +2x_{170}\,x_{172} +x_{171}^{2}
\end{align*}
\begin{align*}
R_{121} &= -x_{136}\,x_{174} -x_{174} +x_{175}
\end{align*}
\begin{align*}
R_{122} &= -1 +x_{173}^{2} +2x_{173}\,x_{175} +x_{174}^{2}
\end{align*}
\begin{align*}
R_{123} &= 1 -x_{136}\,x_{177} -x_{177} +x_{178}
\end{align*}
\begin{align*}
R_{124} &= -1 +2x_{176} +x_{176}^{2} +2x_{176}\,x_{178} +x_{177}^{2}
\end{align*}
\begin{align*}
R_{125} &= -x_{136}\,x_{180} -x_{180} +x_{181}
\end{align*}
\begin{align*}
R_{126} &= -1 +x_{179}^{2} +2x_{179}\,x_{181} +x_{180}^{2}
\end{align*}
\begin{align*}
R_{127} &= -x_{136}\,x_{183} -x_{183} +x_{184}
\end{align*}
\begin{align*}
R_{128} &= -1 +x_{182}^{2} +2x_{182}\,x_{184} +x_{183}^{2}
\end{align*}
\begin{align*}
R_{129} &= -x_{136}\,x_{186} -x_{186} +x_{187}
\end{align*}
\begin{align*}
R_{130} &= -1 +x_{185}^{2} +2x_{185}\,x_{187} +x_{186}^{2}
\end{align*}
\begin{align*}
R_{131} &= -3 -2x_{0} -4x_{0}^{2} -2x_{0}\,x_{3} -x_{1}^{2} -x_{2}^{2} -2x_{3}^{2} -x_{4}^{2} +2x_{8} -x_{8}^{2} +2x_{9}\,x_{10} \\
&\quad {} -x_{11}^{2} -x_{12}^{2} -x_{13}^{2} -x_{14}^{2} -x_{15}^{2} -4x_{16}^{2} -x_{17}^{2} -x_{23} -x_{94}^{2} +2x_{96} \\
&\quad {} -2x_{96}^{2} -x_{97}^{2} -x_{98}^{2} -x_{102}^{2} -x_{105}^{2} -x_{108}^{2} -x_{111}^{2} -x_{114}^{2} -x_{115}^{2} \\
&\quad {} -x_{118}^{2} -x_{121}^{2} -x_{122}^{2} -x_{126}^{2} -x_{129}^{2} -x_{132}^{2} +x_{135} -x_{188}^{2}
\end{align*}
\begin{align*}
R_{132} &= 2x_{4} -4x_{4}\,x_{8} -4x_{4}\,x_{9} -4x_{4}\,x_{10} -4x_{4}\,x_{11} -4x_{4}\,x_{12} -4x_{4}\,x_{13} -4x_{4}\,x_{14} \\
&\quad {} -4x_{4}\,x_{15} -8x_{4}\,x_{16} +x_{188}
\end{align*}
\begin{align*}
R_{133} &= -16 +x_{0} -4x_{189}^{2}
\end{align*}
\begin{align*}
R_{134} &= -x_{136}\,x_{191} -x_{191} +x_{200}
\end{align*}
\begin{align*}
R_{135} &= -1 +x_{190}^{2} +x_{191}^{2} -x_{200}^{2}
\end{align*}
\begin{align*}
R_{136} &= -x_{193} -x_{193}\,x_{194} +x_{201}
\end{align*}
\begin{align*}
R_{137} &= -1 +x_{192}^{2} +x_{193}^{2} -x_{201}^{2}
\end{align*}
\begin{align*}
R_{138} &= -x_{186}^{2} +x_{202}
\end{align*}
\begin{align*}
R_{139} &= x_{191} -2x_{196}\,x_{202} -x_{202}
\end{align*}
\begin{align*}
R_{140} &= -4x_{191} -8x_{191}\,x_{195} +x_{194}
\end{align*}
\begin{align*}
R_{141} &= -x_{136} -x_{190}\,x_{197} +x_{194}
\end{align*}
\begin{align*}
R_{142} &= -x_{185} -x_{187} -x_{190}\,x_{198} +x_{192}
\end{align*}
\begin{align*}
R_{143} &= -x_{94} -4x_{186}\,x_{199} +x_{193}
\end{align*}
\begin{align*}
R_{144} &= 26 +2x_{0}^{2} +2x_{0}\,x_{3} -2x_{0}\,x_{141} -2x_{0}\,x_{150} +x_{1}^{2} -2x_{1}\,x_{177} +x_{2}^{2} -2x_{2}\,x_{168} \\
&\quad {} +x_{3}^{2} -2x_{3}\,x_{150} +x_{8}^{2} -2x_{8}\,x_{153} +x_{9}^{2} -2x_{9}\,x_{156} +x_{10}^{2} -2x_{10}\,x_{159} +x_{11}^{2} \\
&\quad {} -2x_{11}\,x_{162} +x_{12}^{2} -2x_{12}\,x_{165} +x_{13}^{2} -2x_{13}\,x_{171} +x_{14}^{2} -2x_{14}\,x_{174} +x_{15}^{2} \\
&\quad {} -2x_{15}\,x_{180} +4x_{16}^{2} -4x_{16}\,x_{183} +x_{94}^{2} -2x_{94}\,x_{186} -x_{136}\,x_{203} -2x_{140} -2x_{142} \\
&\quad {} +x_{142}^{2} -2x_{149} -2x_{151} +x_{151}^{2} -2x_{152} +2x_{153} -2x_{154} +x_{154}^{2} -2x_{155} \\
&\quad {} -2x_{157} +x_{157}^{2} -2x_{158} -2x_{160} +x_{160}^{2} -2x_{161} -2x_{163} +x_{163}^{2} -2x_{164} \\
&\quad {} -2x_{166} +x_{166}^{2} -2x_{167} -2x_{169} +x_{169}^{2} -2x_{170} -2x_{172} +x_{172}^{2} -2x_{173} \\
&\quad {} -2x_{175} +x_{175}^{2} -2x_{176} +x_{178}^{2} -2x_{179} -2x_{181} +x_{181}^{2} -2x_{182} -2x_{184} \\
&\quad {} +x_{184}^{2} -2x_{185} -2x_{187} +x_{187}^{2}
\end{align*}
\begin{align*}
R_{145} &= -x_{98}\,x_{204} +x_{122}
\end{align*}
\begin{align*}
R_{146} &= -x_{115}\,x_{205} +x_{122}
\end{align*}
\begin{align*}
R_{147} &= -1 +x_{2} -x_{3} -2x_{3}\,x_{212}
\end{align*}
\begin{align*}
R_{148} &= -2 -x_{94}^{2} -x_{141}^{2} -x_{150}^{2} +2x_{153} -x_{153}^{2} -x_{156}^{2} -x_{159}^{2} -x_{162}^{2} -x_{165}^{2} \\
&\quad {} -x_{168}^{2} -x_{171}^{2} -x_{174}^{2} -x_{177}^{2} -x_{180}^{2} -x_{183}^{2} +x_{186}^{2} -x_{206}^{2} -x_{207}^{2} -x_{208} \\
&\quad {} -x_{208}^{2}
\end{align*}
\begin{align*}
R_{149} &= -1 -4x_{100}^{2} -4x_{104}^{2} -4x_{107}^{2} -4x_{110}^{2} -4x_{113}^{2} -4x_{117}^{2} -4x_{120}^{2} \\
&\quad {} +x_{122}^{2} -4x_{124}^{2} -4x_{128}^{2} -4x_{131}^{2} -4x_{134}^{2} -x_{209}^{2} -x_{210}^{2} -x_{211} -x_{211}^{2}
\end{align*}
\begin{align*}
R_{150} &= -x_{3}\,x_{12} +x_{213}
\end{align*}
\begin{align*}
R_{151} &= x_{96}\,x_{215} -2x_{97}\,x_{150} +2x_{102} -2x_{102}\,x_{153} -2x_{105}\,x_{156} -2x_{108}\,x_{159} -2x_{111}\,x_{162} \\
&\quad {} -2x_{114}\,x_{165} -2x_{118}\,x_{171} -2x_{121}\,x_{174} -2x_{126}\,x_{180} -2x_{129}\,x_{183} -2x_{132}\,x_{186} \\
&\quad {} -2x_{136}\,x_{215} +2x_{149}\,x_{150} +2x_{150}\,x_{151} -2x_{152} +2x_{152}\,x_{153} +2x_{153}\,x_{154} -2x_{154} \\
&\quad {} +2x_{155}\,x_{156} +2x_{156}\,x_{157} +2x_{158}\,x_{159} +2x_{159}\,x_{160} +2x_{161}\,x_{162} +2x_{162}\,x_{163} \\
&\quad {} +2x_{164}\,x_{165} +2x_{165}\,x_{166} +2x_{170}\,x_{171} +2x_{171}\,x_{172} +2x_{173}\,x_{174} +2x_{174}\,x_{175} \\
&\quad {} +2x_{179}\,x_{180} +2x_{180}\,x_{181} +2x_{182}\,x_{183} +2x_{183}\,x_{184} +2x_{185}\,x_{186} +2x_{186}\,x_{187} \\
&\quad {} +x_{214} -2x_{215}
\end{align*}
\begin{align*}
R_{152} &= 1 +x_{94}^{2} -x_{96}\,x_{214} +x_{97}^{2} -2x_{97}\,x_{149} +x_{102}^{2} -2x_{102}\,x_{152} +x_{105}^{2} \\
&\quad {} -2x_{105}\,x_{155} +x_{108}^{2} -2x_{108}\,x_{158} +x_{111}^{2} -2x_{111}\,x_{161} +x_{114}^{2} -2x_{114}\,x_{164} \\
&\quad {} +x_{118}^{2} -2x_{118}\,x_{170} +x_{121}^{2} -2x_{121}\,x_{173} +x_{126}^{2} -2x_{126}\,x_{179} +x_{129}^{2} \\
&\quad {} -2x_{129}\,x_{182} +x_{132}^{2} -2x_{132}\,x_{185} +x_{141}^{2} +x_{149}^{2} +x_{152}^{2} +x_{155}^{2} +x_{158}^{2} \\
&\quad {} +x_{161}^{2} +x_{164}^{2} +x_{168}^{2} +x_{170}^{2} +x_{173}^{2} +x_{177}^{2} +x_{179}^{2} +x_{182}^{2} +x_{185}^{2} \\
&\quad {} -2x_{186}^{2} +x_{206}^{2} +x_{207}^{2} +x_{208} +x_{208}^{2} +x_{215}
\end{align*}
\begin{align*}
R_{153} &= x_{96}\,x_{217} -2x_{98}\,x_{141} -2x_{115}\,x_{168} -2x_{122}\,x_{177} -2x_{136}\,x_{217} +2x_{140}\,x_{141} \\
&\quad {} +2x_{141}\,x_{142} +2x_{167}\,x_{168} +2x_{168}\,x_{169} +2x_{176}\,x_{177} +2x_{177} +2x_{177}\,x_{178} +x_{216} \\
&\quad {} -3x_{217}
\end{align*}
\begin{align*}
R_{154} &= -x_{96}\,x_{216} +x_{98}^{2} -2x_{98}\,x_{140} +x_{115}^{2} -2x_{115}\,x_{167} +x_{122}^{2} -2x_{122}\,x_{176} +x_{140}^{2} \\
&\quad {} -x_{141}^{2} +x_{167}^{2} -x_{168}^{2} +x_{176}^{2} -x_{177}^{2} +x_{216} +x_{217}
\end{align*}

\endgroup

\section{Artifact identities}
\par\smallskip\noindent\begin{minipage}{\linewidth}
\textbf{Pinned 242-coordinate source.} SHA-256:
\begin{quote}\small\ttfamily
\nolinkurl{50f88c245d3a32b5c041924ea4ac07b9149e5da24c859f3a167ea984d93e18d2}
\end{quote}
\end{minipage}\par
\par\smallskip\noindent\begin{minipage}{\linewidth}
\textbf{Composed 223-coordinate parent.} SHA-256:
\begin{quote}\small\ttfamily
\nolinkurl{ebf82c1646ff38125aaed2bdef0e7fc5c07b67853dda0e781fc459ead0d1e7c4}
\end{quote}
\end{minipage}\par
\par\smallskip\noindent\begin{minipage}{\linewidth}
\textbf{Principal residual system.} SHA-256:
\begin{quote}\small\ttfamily
\nolinkurl{95c0db6befd4ad98e6749ba18580f3dd775b1faaec36d270f8b839911754d6b1}
\end{quote}
\end{minipage}\par
\par\smallskip\noindent\begin{minipage}{\linewidth}
\textbf{Companion residual system.} SHA-256:
\begin{quote}\small\ttfamily
\nolinkurl{a91308bd0c82182ea15080763cb23cb4fa0aa45a812eabb7a6cbcd21c0b95a06}
\end{quote}
\end{minipage}\par
\par\smallskip\noindent\begin{minipage}{\linewidth}
\textbf{Principal collected quartic.} SHA-256:
\begin{quote}\small\ttfamily
\nolinkurl{a3ee776477f2b6ac5eca2e5e55726e916f596317170d4d92452a2363504c400a}
\end{quote}
\end{minipage}\par
\par\smallskip\noindent\begin{minipage}{\linewidth}
\textbf{Companion collected quartic.} SHA-256:
\begin{quote}\small\ttfamily
\nolinkurl{e5f7a57354358df48a06d9b82ff1bafc3baecde8ce927b1a504c7c015b79ccf9}
\end{quote}
\end{minipage}\par


The algebra certificates contain the five radical directions,
both affine combinations, the centered chart and its inverse,
and every row operation. The source-composition certificates
contain the cofinal-range edits, private-wire maps, pooled
blocks, and complete frame-atom ledger.

\begin{thebibliography}{9}
\bibitem{s157}
ProveIt contributors.
\emph{Affine beta spacing and arithmetic sign recovery;
source 157, linear CRT}.
Snapshot \pp{953fb4d82d4229f4efc5f9db4cc2bf38ae329bb4},
25 October 2026.
\url{https://github.com/VladimirReshetnikov/ProveIt/tree/953fb4d82d4229f4efc5f9db4cc2bf38ae329bb4/Computability/DiophantineConsistency/con-zfc-diophantine-forms/157-linear-crt}.

\bibitem{s158}
ProveIt contributors.
\emph{Pooling exact power certificates for Diophantine consistency;
source 158}.
Same pinned snapshot.
\url{https://github.com/VladimirReshetnikov/ProveIt/tree/953fb4d82d4229f4efc5f9db4cc2bf38ae329bb4/Computability/DiophantineConsistency/con-zfc-diophantine-forms/158-pooled-powers}.

\bibitem{s159}
ProveIt contributors.
\emph{A direct quartic for Con(ZFC) in 234 integer variables;
source 159}.
Same pinned snapshot.
\url{https://github.com/VladimirReshetnikov/ProveIt/tree/953fb4d82d4229f4efc5f9db4cc2bf38ae329bb4/Computability/DiophantineConsistency/con-zfc-diophantine-forms/159-quadratic-projection-234}.

\bibitem{endpoint}
ProveIt contributors.
\emph{KEndpoint.lean}, especially
\pp{c222_iff_zf_refutation}.
\url{https://github.com/VladimirReshetnikov/ProveIt/blob/953fb4d82d4229f4efc5f9db4cc2bf38ae329bb4/Computability/BusyBeaver/Lean/BusyBeaver/ZFC/KEndpoint.lean}.

\bibitem{choice}
ProveIt contributors.
\emph{ConstructibleChoice.lean}, especially
\pp{zf_refutation_iff_zfc_refutation}.
\url{https://github.com/VladimirReshetnikov/ProveIt/blob/953fb4d82d4229f4efc5f9db4cc2bf38ae329bb4/Computability/BusyBeaver/Lean/BusyBeaver/ZFC/ConstructibleChoice.lean}.

\bibitem{current}
ProveIt contributors.
\emph{A complete fifteen-variable consistency equation}.
Complete display at the same snapshot.
\url{https://github.com/VladimirReshetnikov/ProveIt/blob/953fb4d82d4229f4efc5f9db4cc2bf38ae329bb4/Computability/DiophantineConsistency/equation.tex}.

\bibitem{late}
ProveIt contributors.
\emph{Sources 160--165: later direct quartic and compact tradeoffs}.
README comparison at commit
\pp{fb1bed5eb7196cf35abfaa746c18c6ee7eae6cfb}.
\url{https://github.com/VladimirReshetnikov/ProveIt/tree/fb1bed5eb7196cf35abfaa746c18c6ee7eae6cfb/Computability/DiophantineConsistency/con-zfc-diophantine-forms}.
\end{thebibliography}
\end{document}
